\section{Results}
\label{sec:results}

\pgfplotsset{xo/.style={grid=major, grid style={gray!18}, tick label style={font=\footnotesize},
  label style={font=\footnotesize}, title style={font=\footnotesize\bfseries}, legend style={font=\footnotesize, draw=none},
  unbounded coords=jump, every axis plot/.append style={line width=0.9pt}}}
\newcommand{\xocurve}[4]{% #1 colour, #2 mark, #3 file, #4 column prefix
  \addplot[#1, mark=#2, mark size=1.4pt, error bars/.cd, y dir=both, y explicit]
    table[x=rate, y=#4p99, y error plus expr=\thisrow{#4hi}-\thisrow{#4p99}, y error minus expr=\thisrow{#4p99}-\thisrow{#4lo}] {#3};}

The campaign comprised \xo{cells} cells (\xo{reps} repetitions of each configuration), ran for \xo{hours}\,h,
and needed no exclusions: no pod restarted and the host was never suspended during a step. All
\xo{equiv/total} equivalence responses were byte-identical to the monolith's, and the \texttt{calls} field
confirmed the hop counts of Table~\ref{tab:hops} in every request. At the lowest rates the generator's sends lagged
the constant-rate schedule by up to \xo{lagmax/browse}\,ms (browse) and \xo{lagmax/checkout}\,ms (checkout) at the
99th percentile, the same order as the latencies being measured. Latency is measured from the actual send, so the
lag does not enter it, but low-rate arrivals were burstier than a constant rate.

\subsection{Latency at low load (RQ1, H1)}
\label{sec:res-latency}

Figure~\ref{fig:open} shows p99 latency against offered load. At the lowest rate every decomposed configuration
is slower than the monolith. For browse, p99 was \xo{p99/browse/A/50}\,ms in~A, \xo{p99/browse/B/50}\,ms in~B,
\xo{p99/browse/C/50}\,ms in~C and \xo{p99/browse/D/50}\,ms in~D; for checkout, \xo{p99/checkout/A/25},
\xo{p99/checkout/B/25}, \xo{p99/checkout/C/25} and \xo{p99/checkout/D/25}\,ms. Table~\ref{tab:h1} gives the
ratios to the monolith with their Bonferroni-adjusted intervals. All six intervals lie above one, so
\textbf{H1 is supported}. Pooling the requests of all repetitions, the p99 of B at low load is \xo{ratio/browse/B/50/est} (browse) and
\xo{ratio/checkout/B/25/est} (checkout) times the monolith's, that of C \xo{ratio/browse/C/50/est} and
\xo{ratio/checkout/C/25/est} times, and that of D \xo{ratio/browse/D/50/est} and \xo{ratio/checkout/D/25/est}
times.

Expressed per hop, the median latency of B and C grew by \xo{perhop/browse/C/p50}--\xo{perhop/checkout/B/p50}\,ms
for every remote call (browse in C and checkout in B, respectively), and that of D by \xo{perhop/browse/D/p50}--%
\xo{perhop/checkout/D/p50}\,ms.

All four configurations were slower at the lowest offered rates than at moderate ones; the monolith's browse p99
fell from \xo{p99/browse/A/50}\,ms at 50 requests per second to \xo{p99/browse/A/800}\,ms at 800. A control run
traced this to the host's idle states rather than to warm-up (Section~\ref{sec:threats}). It inflates absolute
low-load latencies of every configuration; the ratios in Table~\ref{tab:h1} compare configurations under the same
condition.

As the offered rate rises, the configurations part ways. The monolith sustained every rate of the grid, up to
\xo{gridtop/browse} requests per second for browse and \xo{gridtop/checkout} for checkout; from 500 requests per
second upwards its p99 never exceeded \xo{p99/checkout/A/1800}\,ms. Configuration~B sustained every level up to \xo{sustained/browse/B} (browse) and \xo{sustained/checkout/B}
(checkout) requests per second in every repetition, C up to \xo{sustained/browse/C} and \xo{sustained/checkout/C},
and D up to \xo{sustained/browse/D} and \xo{sustained/checkout/D}. Near their limits the decomposed configurations' tails become erratic. At 500
browse requests per second, C's median p99 across repetitions was \xo{p99/browse/C/500}\,ms, but the repetitions
ranged from \xo{p99lo/browse/C/500} to \xo{p99hi/browse/C/500}\,ms. The edge service, \texttt{search}, was then
using \xo{hot/browse/C/500/util}\% of its half-core limit on average, and \xo{hot/browse/C/500/thr}\% of its CFS
periods were throttled. This is the mechanism of Section~\ref{sec:bg-cfs}: well before the limit is reached on
average, bursts exhaust the quota within a period and the requests that arrive next wait for the following one.
B shows the same pattern beyond its sustained rate (\texttt{catalog} at \xo{hot/browse/B/2200/util}\% with
\xo{hot/browse/B/2200/thr}\% of periods throttled at 2{,}200 requests per second, p99 \xo{p99/browse/B/2200}\,ms).
The monolith's application container used at most \xo{hot/browse/A/3200/util}\% of its four-core limit in the
open-loop sweeps and was never throttled.


\begin{figure}[t]
\centering
\begin{tikzpicture}
\begin{groupplot}[group style={group size=2 by 1, horizontal sep=1.5cm}, xo, width=0.5\textwidth, height=5.4cm,
    xmode=log, ymode=log, xlabel={Offered load (requests/s)}, ylabel={p99 latency (ms)}]
\nextgroupplot[title={Browse}, legend columns=4, legend style={at={(1.1,1.2)}, anchor=south}]
\xocurve{xoA}{*}{\gen{open-browse.csv}}{A}\addlegendentry{A (monolith)}
\xocurve{xoB}{square*}{\gen{open-browse.csv}}{B}\addlegendentry{B (4 services)}
\xocurve{xoC}{triangle*}{\gen{open-browse.csv}}{C}\addlegendentry{C (8 services)}
\xocurve{xoD}{diamond*}{\gen{open-browse.csv}}{D}\addlegendentry{D (8 + Istio)}
\nextgroupplot[title={Checkout}]
\xocurve{xoA}{*}{\gen{open-checkout.csv}}{A}
\xocurve{xoB}{square*}{\gen{open-checkout.csv}}{B}
\xocurve{xoC}{triangle*}{\gen{open-checkout.csv}}{C}
\xocurve{xoD}{diamond*}{\gen{open-checkout.csv}}{D}
\end{groupplot}
\end{tikzpicture}
\caption{Tail latency under open-loop load: median p99 across \xo{reps} repetitions, bars spanning the
repetitions' spread. A line ends where the configuration saturated (more than 5\% errors, throughput below 90\% of the offered rate, or
p99 above 100\,ms). An isolated mark is a level sustained again after one that was not (C at 600 browse
requests per second).}
\label{fig:open}
\end{figure}

\begin{table}[t]
\caption{H1: ratio of each microservice configuration's p99 to the monolith's at the lowest offered rate, with
two-level bootstrap intervals at the Bonferroni-adjusted level $1-0.05/6$. A ratio above one means the monolith
is faster.}
\label{tab:h1}
\centering\small
\begin{tabular}{@{}lrlrl@{}}
\toprule
Workload & Rate (req/s) & Configuration & p99 ratio & Interval \\
\midrule
\xoinput{\gen{tab-h1.tex}}
\bottomrule
\end{tabular}
\end{table}

\subsection{Capacity under an equal budget (RQ2, H2)}
\label{sec:res-capacity}

Figure~\ref{fig:closed} shows closed-loop throughput as a function of the number of concurrent clients, and
Table~\ref{tab:usl} the fitted USL parameters. The monolith reached a peak of \xo{usl/browse/A/xmax} browse and
\xo{usl/checkout/A/xmax} checkout requests per second. The decomposed configurations peaked far lower: B at
\xo{xmaxratio/browse/B} (browse) and \xo{xmaxratio/checkout/B} (checkout) of the monolith's capacity, C at
\xo{xmaxratio/browse/C} and \xo{xmaxratio/checkout/C}, and D at \xo{xmaxratio/browse/D} and
\xo{xmaxratio/checkout/D}. No decomposed configuration matched the monolith at any measured concurrency level in
either workload, so the throughput crossover does not exist within the tested range and \textbf{H2 is falsified} by
its first condition. Its second condition is met only formally: the fits for B and C put $\beta$ at its lower
bound of zero, which is lower than the monolith's \xo{usl/browse/A/beta} (browse), but for the reason discussed
next rather than because decomposition removed a coherency cost.

The cgroup counters show what limits each configuration. At 64 clients the monolith's database container used
\xo{chot/browse/A/64/util}\% of its two-core limit and was throttled in \xo{chot/browse/A/64/thr}\% of its periods,
while the application container used \xo{chot/browse/A/64/apputil}\% of its four cores (checkout:
\xo{chot/checkout/A/64/util}\% and \xo{chot/checkout/A/64/apputil}\%). The monolith is therefore bound by its
database tier, and its fitted $\beta>0$ and $N^\ast$ of \xo{usl/browse/A/nstar} (browse) and \xo{usl/checkout/A/nstar}
(checkout) describe a saturated database rather than application code. In every decomposed configuration the
limit is instead the CPU quota of the edge service: \texttt{\xo{chot/browse/B/64/appname}} in~B at
\xo{chot/browse/B/64/apputil}\% of its single core and \texttt{\xo{chot/browse/C/64/appname}} in~C at
\xo{chot/browse/C/64/apputil}\% of its half core, while their databases stayed below saturation.
The edge service does the most work per request: it parses its callees' JSON responses, encodes the client's
response, and drives every outbound call, so it is the first to exhaust an equal share of the budget.

Configuration~C is the extreme case. Its USL fit has $\alpha = \xo{usl/browse/C/alpha}$ and the throughput curve is
flat from a single client onwards (\xo{x/browse/C/1} requests per second at $N=1$, \xo{x/browse/C/256} at $N=256$).
With a single client, \texttt{search} already used \xo{chot/browse/C/1/apputil}\% of its half-core limit and was
throttled in \xo{chot/browse/C/1/thr}\% of its CFS periods; additional clients only queue. In USL terms the configuration is completely contention-bound, and the
contended resource is a CPU quota. B and D have $\alpha$ between \xo{usl/checkout/D/alpha} and \xo{usl/browse/B/alpha}. D's throughput with one
client is lower than C's because of the sidecars' added latency (\xo{x/browse/D/1} against \xo{x/browse/C/1}
requests per second), but its peak is slightly higher (\xo{usl/browse/D/xmax} against \xo{usl/browse/C/xmax}). Two
measured effects explain this: the sidecars run on CPU outside the application budget, and D's edge application
container needed less CPU per request than C's (\xo{hot/browse/D/500/util}\% against \xo{hot/browse/C/500/util}\%
of the same limit at 500 requests per second).


\begin{figure}[t]
\centering
\begin{tikzpicture}
\begin{groupplot}[group style={group size=2 by 1, horizontal sep=1.5cm}, xo, width=0.5\textwidth, height=5.4cm,
    xmode=log, ymode=log, log basis x=2, xlabel={Concurrent clients $N$}, ylabel={Throughput (requests/s)},
    xtick={1,4,16,64,256}, xticklabels={1,4,16,64,256}]
\nextgroupplot[title={Browse}, legend columns=4, legend style={at={(1.1,1.2)}, anchor=south}]
\foreach \c/\m in {A/*, B/square*, C/triangle*, D/diamond*} {
  \edef\tmp{\noexpand\addplot[xo\c, only marks, mark=\m, mark size=1.5pt, forget plot, error bars/.cd, y dir=both, y explicit] table[x=n, y=\c x, y error plus expr=\noexpand\thisrow{\c hi}-\noexpand\thisrow{\c x}, y error minus expr=\noexpand\thisrow{\c x}-\noexpand\thisrow{\c lo}] {\gen{closed-browse.csv}};}\tmp
  \edef\tmp{\noexpand\addplot[xo\c, no marks] table[x=n, y=\c fit] {\gen{usl-browse.csv}};}\tmp
}
\legend{A (monolith), B (4 services), C (8 services), D (8 + Istio)}
\nextgroupplot[title={Checkout}]
\foreach \c/\m in {A/*, B/square*, C/triangle*, D/diamond*} {
  \edef\tmp{\noexpand\addplot[xo\c, only marks, mark=\m, mark size=1.5pt, forget plot, error bars/.cd, y dir=both, y explicit] table[x=n, y=\c x, y error plus expr=\noexpand\thisrow{\c hi}-\noexpand\thisrow{\c x}, y error minus expr=\noexpand\thisrow{\c x}-\noexpand\thisrow{\c lo}] {\gen{closed-checkout.csv}};}\tmp
  \edef\tmp{\noexpand\addplot[xo\c, no marks] table[x=n, y=\c fit] {\gen{usl-checkout.csv}};}\tmp
}
\end{groupplot}
\end{tikzpicture}
\caption{Closed-loop throughput (marks: median across \xo{reps} repetitions, bars spanning their range) and
fitted Universal Scalability Law curves (lines). Both axes are logarithmic.}
\label{fig:closed}
\end{figure}

\begin{table}[t]
\caption{Capacity under the equal budget. Peak: highest median closed-loop throughput (requests/s) over the
concurrency levels, at $N$ clients, with a 95\% interval from resampling repetitions. Bound: operational bottleneck
bound $\min_k X/U_k$ at the highest sustained open-loop rate, with the container that attains it. Budget used: CPU
of all application and database containers at 64 clients, as a share of the 6-core budget.}
\label{tab:capacity}
\centering\scriptsize
\setlength{\tabcolsep}{3pt}
\begin{tabular}{@{}lllrrlr@{}}
\toprule
Workload & Configuration & Peak [95\% interval] & $N$ & vs.\ A & Bound (container) & Budget used \\
\midrule
\xoinput{\gen{tab-capacity.tex}}
\bottomrule
\end{tabular}
\end{table}

\begin{table}[t]
\caption{USL fits to the closed-loop measurements, a secondary description of the curves in
Figure~\ref{fig:closed}. Brackets: 95\% intervals from refits on resampled repetitions. The fit constrains
$0 \le \alpha, \beta \le 1$; $\beta = 0$ means the fit sits at that bound, so the curve has no retrograde region,
$N^\ast$ is undefined and $X_{\max}$ is $\lambda/\alpha$.}
\label{tab:usl}
\centering\scriptsize
\setlength{\tabcolsep}{3pt}
\begin{tabular}{@{}llrllrrr@{}}
\toprule
Workload & Configuration & $\lambda$ & $\alpha$ & $\beta$ & $N^\ast$ & $X_{\max}$ & $R^2_{\mathrm{adj}}$ \\
\midrule
\xoinput{\gen{tab-usl.tex}}
\bottomrule
\end{tabular}
\end{table}

The USL describes the monolith well ($R^2_{\mathrm{adj}}$ of \xo{usl/browse/A/r2adj} and \xo{usl/checkout/A/r2adj})
and B and D only moderately (\xo{usl/browse/D/r2adj}--\xo{usl/checkout/D/r2adj}): their throughput climbs to the
quota within a few clients and then stays flat, a shape between the model's two regimes. For C the curve is flat
from the start, and $R^2$ carries no information.

\subsection{Where the time goes (RQ3)}
\label{sec:res-layers}

Figure~\ref{fig:layers} breaks the mean latency at the lowest rate into the six layers. In the monolith most
of it is database time (\xo{layer/browse/A/50/db} of \xo{mean/browse/A/50}\,ms for browse,
\xo{layer/checkout/A/25/db} of \xo{mean/checkout/A/25}\,ms for checkout); application code, serialisation and
framework dispatch together take a few tenths of a millisecond. Decomposition leaves database time roughly where
it was (browse: \xo{layer/browse/B/50/db}\,ms in~B and \xo{layer/browse/C/50/db}\,ms in~C) and adds inter-service
transit on top: \xo{layer/browse/B/50/net} and \xo{layer/browse/C/50/net}\,ms for browse in B and~C,
\xo{layer/checkout/B/25/net} and \xo{layer/checkout/C/25/net}\,ms for checkout. Per hop this is
\xo{perhop/browse/C/net}--\xo{perhop/browse/B/net}\,ms for browse and \xo{perhop/checkout/C/net}--%
\xo{perhop/checkout/B/net}\,ms for checkout. Serialisation, often blamed for the cost of microservices, adds only
\xo{perhop/browse/D/ser}--\xo{perhop/checkout/B/ser}\,ms per hop for these payloads (at most about 5.5\,kB of JSON).
In D the transit layer nearly doubles again and database time rises, because every query now also passes
through two sidecars (Section~\ref{sec:res-mesh}).

The decomposed configurations also cost more to run. Table~\ref{tab:cost} reports CPU time per request at the
highest rate all four sustained. For browse at \xo{ref/browse/high} requests per second the monolith spent
\xo{cpums/browse/A/600/total}\,ms of CPU per request, B \xo{cpums/browse/B/600/total}, C
\xo{cpums/browse/C/600/total} and D \xo{cpums/browse/D/600/total}\,ms; for checkout at \xo{ref/checkout/high} the
figures are \xo{cpums/checkout/A/400/total}, \xo{cpums/checkout/B/400/total}, \xo{cpums/checkout/C/400/total} and
\xo{cpums/checkout/D/400/total}\,ms. Each remote call added \xo{perhop/browse/B/cpu}--\xo{perhop/browse/C/cpu}\,ms of
CPU in B and~C and \xo{perhop/browse/D/cpu}--\xo{perhop/checkout/D/cpu}\,ms in~D. The resident memory of all pods
grew from \xo{mem/browse/A/600}\,GiB for the monolith to \xo{mem/browse/B/600}, \xo{mem/browse/C/600} and
\xo{mem/browse/D/600}\,GiB, mostly the fixed footprint of one JVM per service. The database tier's CPU per request
rose more slowly, from \xo{cpums/browse/A/600/db} to \xo{cpums/browse/D/600/db}\,ms for browse; most of the extra
cost is in the application and mesh tiers.


\begin{figure}[t]
\centering
\begin{tikzpicture}
\begin{groupplot}[group style={group size=2 by 1, horizontal sep=0.6cm, yticklabels at=edge left}, xo, width=0.52\textwidth,
    height=4.6cm, xbar stacked, /pgf/bar width=7pt, xmin=0, ytick={0,1,2,3}, yticklabels={A,B,C,D}, y dir=reverse,
    xlabel={Mean latency (ms)}, enlarge y limits=0.18,
    legend image code/.code={\fill[#1] (0cm,-0.08cm) rectangle (0.22cm,0.08cm);}]
\nextgroupplot[title={Browse, \xo{ref/browse/low} req/s}, legend columns=6, legend style={at={(1.03,1.24)}, anchor=south, /tikz/every even column/.append style={column sep=4pt}}]
\foreach \l/\c in {app/layApp, db/layDb, ser/laySer, net/layNet, oth/layOth, edge/layEdge} {
  \edef\tmp{\noexpand\addplot[fill=\c, draw=white, line width=0.6pt] table[x=\l, y=cfg] {\gen{layers-browse.csv}};}\tmp
}
\legend{Application, Database, Serialisation, Inter-service network, Framework, Edge}
\nextgroupplot[title={Checkout, \xo{ref/checkout/low} req/s}]
\foreach \l/\c in {app/layApp, db/layDb, ser/laySer, net/layNet, oth/layOth, edge/layEdge} {
  \edef\tmp{\noexpand\addplot[fill=\c, draw=white, line width=0.6pt] table[x=\l, y=cfg] {\gen{layers-checkout.csv}};}\tmp
}
\end{groupplot}
\end{tikzpicture}
\caption{Where the time goes at the lowest offered rate: mean per-request latency by layer (median across
repetitions), from the \texttt{Server-Timing} accounting of Section~\ref{sec:timing}.}
\label{fig:layers}
\end{figure}

\begin{table}[t]
\caption{Resource cost at the highest offered rate every configuration sustained (browse \xo{ref/browse/high}
and checkout \xo{ref/checkout/high} requests/s): CPU milliseconds per request by tier (cgroup
\texttt{cpu.stat}), resident memory of all pods, the busiest application container with its CPU utilisation
(share of its limit) and the share of CFS periods in which it was throttled, and the resulting p99.}
\label{tab:cost}
\centering\small
\setlength{\tabcolsep}{3.5pt}
\begin{tabular}{@{}llrrrrlrrr@{}}
\toprule
& & \multicolumn{3}{c}{CPU ms per request} & Memory & \multicolumn{3}{c}{Busiest application container} & p99 \\
\cmidrule(lr){3-5}\cmidrule(lr){7-9}
Workload & Configuration & app & db & mesh & (GiB) & service & util. & throttled & (ms) \\
\midrule
\xoinput{\gen{tab-cost.tex}}
\bottomrule
\end{tabular}
\end{table}

\subsection{The cost of the service mesh (RQ4, H3)}
\label{sec:res-mesh}

Table~\ref{tab:mesh} lists the mesh's cost at every rate where the comparison is clean. Per hop, the two
sidecars added \xo{mesh/browse/200/perhop}--\xo{mesh/checkout/25/perhop}\,ms of transit, of the same order as
the 0.18--0.25\,ms (P90--P99) that Istio reports for a sidecar pair in version 1.22~\cite{Istio122Perf}, and the database
round trips of a request took \xo{mesh/browse/200/db}--\xo{mesh/checkout/200/db}\,ms longer. The mesh explained
\xo{mesh/browse/50/share}\% of D's mean end-to-end latency on browse at the lowest rate and
\xo{mesh/checkout/25/share}\% on checkout. The share rose with load, to \xo{mesh/browse/400/share}\% and
\xo{mesh/checkout/200/share}\%. The difference D$-$C itself stayed nearly constant; the share rose because D's own
latency fell with load, as the host's low-load penalty receded.

H3 requires the mesh to account for at least a quarter of end-to-end latency on the six-hop path and its per-hop
cost not to shrink as the path gets deeper. On checkout at the lowest rate, where the condition is evaluated, the
mesh accounted for \xo{H3/share_deep}\% of end-to-end latency (median over all clean rates:
\xo{H3/share_deep_median}\%). At each workload's lowest rate the per-hop transit cost was
\xo{H3/per_hop_shallow_ms}\,ms on the three-hop browse path and \xo{H3/per_hop_deep_ms}\,ms on the six-hop
checkout path. It did not shrink as the path got longer, so the total grows at least linearly with the number of
hops. The two
per-hop intervals overlap, so the data show linear growth, not compounding. \textbf{H3 is supported.} Our original
statement of H3 also predicted that the mesh would become the single largest contributor above four hops. On
checkout at the lowest rate the mesh's \xo{mesh/checkout/25/e2e}\,ms is indeed larger than any single layer of
configuration~C (whose largest, database time, is \xo{layer/checkout/C/25/db}\,ms).

The sidecars' CPU is the larger cost. At \xo{ref/browse/high} browse requests per second they consumed
\xo{cpumsref/browse/D/mesh}\,ms of CPU per request, more than the whole application tier
(\xo{cpumsref/browse/D/app}\,ms); for checkout at \xo{ref/checkout/high} requests per second,
\xo{cpumsref/checkout/D/mesh} against \xo{cpumsref/checkout/D/app}\,ms. Because the sidecars' CPU lies outside the
equal application budget, D's capacity is not lower than C's; under a budget that included them, it would very
likely be lower.


\begin{table}[t]
\caption{Cost of the Istio mesh, measured as D minus C at equal offered rate (medians across repetitions; brackets: 95\% bootstrap interval of the difference of medians, per hop). Share is the part of D's mean end-to-end latency
that the difference explains.}
\label{tab:mesh}
\centering\small
\begin{tabular}{@{}lrrlrrr@{}}
\toprule
Workload & Rate (req/s) & Hops & Transit per hop (ms) & DB (ms) & End to end (ms) & Share \\
\midrule
\xoinput{\gen{tab-mesh.tex}}
\bottomrule
\end{tabular}
\end{table}

\subsection{Capacity under demand-proportional allocation (RQ2, H2)}
\label{sec:res-alloc}

\begin{table}[t]
\caption{Equal versus demand-proportional CPU allocation in the local environment: measured peak throughput
(requests/s), the same relative to the monolith measured in the same campaign, concurrency levels at
which a configuration measured above the monolith, and the busiest container at 64 clients with its CPU
utilisation.}
\label{tab:alloc}
\centering\scriptsize
\setlength{\tabcolsep}{3pt}
\begin{tabular}{@{}llrrrrlll@{}}
\toprule
& & \multicolumn{2}{c}{Peak} & \multicolumn{2}{c}{vs.\ A} & & \multicolumn{2}{c}{Busiest container at $N=64$} \\
\cmidrule(lr){3-4}\cmidrule(lr){5-6}\cmidrule(l){8-9}
Workload & Configuration & equal & prop. & equal & prop. & Wins over A & equal & prop. \\
\midrule
\xoinput{generated/tab-alloc.tex}
\bottomrule
\end{tabular}
\end{table}

Table~\ref{tab:alloc} compares the equal split with the demand-proportional allocation of
Section~\ref{sec:treatments}. Relative to the monolith measured alongside, moving the budget to where the demand
was raised the capacity of the decomposed configurations by a factor of \xo{alloc/gainmin} to
\xo{alloc/gainmax}. B reached \xo{prop/peakpct/browse/B}\% of the monolith's peak on browse and
\xo{prop/peakpct/checkout/B}\% on checkout, against \xo{peakpct/browse/B}\% and \xo{peakpct/checkout/B}\% under
the equal split. C reached \xo{prop/peakpct/browse/C}\% and \xo{prop/peakpct/checkout/C}\%, against
\xo{peakpct/browse/C}\% and \xo{peakpct/checkout/C}\%. The monolith's allocation is unchanged, and its peak of
\xo{prop/peak/browse/A} browse requests per second was close to its \xo{peak/browse/A} in the first
campaign. No decomposed configuration matched the monolith at any measured concurrency level, so \textbf{H2 is
falsified under this treatment as well}.

The busiest containers explain why. The monolith was still bound by its database. B and C were still bound by
their edge services, now at the limit of a two- to fourfold larger quota (\texttt{\xo{alloc/edge/browse/C}} in~C:
\xo{alloc/cores/browse/C} cores instead of \xo{alloc/equal/C}). The exception is checkout in~B, where
\texttt{order} was no longer the limit and \texttt{inventory} took its place. Reallocation removed the idle shares
but not the extra work. At 600 browse requests per second, C's application tier needed
\xo{prop/cpums/browse/C/600/app}\,ms of CPU per request against the monolith's \xo{prop/cpums/browse/A/600/app}\,ms.
At that cost even a perfect split of the four application cores serves fewer than half the requests per second
at which the monolith's database saturated.

H1 and H3 remained supported. The p99 ratios to the monolith at the lowest rate were
\xo{prop/h1/min}--\xo{prop/h1/max}, against
\xo{h1/min}--\xo{h1/max} under the equal split. The allocation did not cause this
difference. At these rates no edge container was throttled in either campaign, and the decomposed configurations'
p99 barely changed (browse in~B: \xo{p99/browse/B/50} and \xo{prop/p99/browse/B/50}\,ms). The monolith, whose
allocation is identical in both campaigns, measured slower in the second one: \xo{p99/browse/A/50} against
\xo{prop/p99/browse/A/50}\,ms on browse and \xo{p99/checkout/A/25} against \xo{prop/p99/checkout/A/25}\,ms on
checkout. Low-load tail latency therefore drifts between campaigns on this host by more than the intervals within
one campaign suggest. Section~\ref{sec:threats} returns to this point.


\subsection{Capacity when the business logic is heavier (RQ2, H2)}
\label{sec:res-work}

\begin{table}[t]
\caption{Service-demand treatment: every module operation spends 1\,ms of CPU (local environment, Configurations A--C,
\xo{work/reps} repetitions). Columns as in Table~\ref{tab:alloc}: measured peak throughput (requests/s) under the
equal split and under demand-proportional allocation, the same relative to the monolith, concurrency levels at
which a configuration measured above the monolith, and the busiest container at 64 clients.}
\label{tab:work}
\centering\scriptsize
\setlength{\tabcolsep}{3pt}
\begin{tabular}{@{}llrrrrlll@{}}
\toprule
& & \multicolumn{2}{c}{Peak} & \multicolumn{2}{c}{vs.\ A} & & \multicolumn{2}{c}{Busiest container at $N=64$} \\
\cmidrule(lr){3-4}\cmidrule(lr){5-6}\cmidrule(l){8-9}
Workload & Configuration & equal & prop. & equal & prop. & Wins over A & equal & prop. \\
\midrule
\xoinput{generated/tab-work.tex}
\bottomrule
\end{tabular}
\end{table}

With 1\,ms of CPU in every module operation the monolith is no longer bound by its database. At 64 clients its
application container used \xo{work/chot/browse/A/64/apputil}\% of its four cores on both workloads, and it peaked
at \xo{work/peak/browse/A} browse and \xo{work/peak/checkout/A} checkout requests per second (Table~\ref{tab:work}).
The gap to the decomposed configurations narrowed. Under the equal split, B reached \xo{work/peakpct/browse/B}\% of
the monolith's peak on browse and \xo{work/peakpct/checkout/B}\% on checkout, and C \xo{work/peakpct/browse/C}\% and
\xo{work/peakpct/checkout/C}\%, against \xo{peakpct/browse/B}--\xo{peakpct/checkout/B}\% and
\xo{peakpct/browse/C}--\xo{peakpct/checkout/C}\% without the added work. With proportional allocation B reached
\xo{workprop/peakpct/browse/B}\% and \xo{workprop/peakpct/checkout/B}\%, and C \xo{workprop/peakpct/browse/C}\% and
\xo{workprop/peakpct/checkout/C}\%. Still no decomposed configuration matched the monolith at any measured
concurrency, so \textbf{H2 is falsified under this treatment as well}.

What remains is the cost of distribution, which is now small relative to the work. At \xo{workprop/ref/browse/high}
browse requests per second C's application tier needed \xo{workprop/cpumsref/browse/C/app}\,ms of CPU per request,
\xo{workprop/cpumsref/browse/A/app}\,ms for the monolith. Even a perfect allocation of the same four cores could therefore serve C at most that fraction of the
monolith's rate, and the measured peaks lie below this bound. The latency penalty
at low load shrank in the same way: the p99 ratios to the monolith were \xo{work/h1/min}--\xo{work/h1/max} under the
equal split, against \xo{h1/min}--\xo{h1/max} without the added work, and H1 remained supported.


% Cloud environment and autoscaling treatment. Not \input yet: the text is written once the EKS and
% EKS-HPA campaigns are analysed (analysis/analyze.py --env eks, compare_env.py, analyze_hpa.py --env eks-hpa).
\subsection{The cloud environment}
\label{sec:res-cloud}

\begin{table}[t]
\caption{Local and cloud environment side by side, equal CPU budget. p99 at the lowest offered rate (browse 50,
checkout 25 requests/s) and its ratio to the monolith (H1); USL peak throughput $X_{\max}$ and its ratio to the
monolith (H2); inter-service transit per hop at low load. Medians across repetitions.}
\label{tab:env}
\centering\scriptsize
\setlength{\tabcolsep}{3pt}
\begin{tabular}{@{}llrrrrrrrrrr@{}}
\toprule
& & \multicolumn{5}{c}{Local (kind)} & \multicolumn{5}{c}{Cloud (EKS)} \\
\cmidrule(lr){3-7}\cmidrule(l){8-12}
Workload & Configuration & p99 & vs.\ A & $X_{\max}$ & vs.\ A & hop & p99 & vs.\ A & $X_{\max}$ & vs.\ A & hop \\
\midrule
\xoinput{generated/tab-env.tex}
\bottomrule
\end{tabular}
\end{table}

% TODO(results): text.

\subsection{Autoscaling under a shared ceiling (RQ2, H2)}
\label{sec:res-hpa}

\begin{figure}[t]
\centering
\begin{tikzpicture}
\begin{groupplot}[group style={group size=2 by 2, horizontal sep=1.5cm, vertical sep=1.6cm}, xo, width=0.5\textwidth,
    height=4.6cm, xmode=log, xlabel={Offered load (requests/s)}]
\nextgroupplot[title={Browse}, ymode=log, xmin=350, xmax=11000, xtick={500,1000,2000,4000,8000}, xticklabels={500,1k,2k,4k,8k}, ylabel={p99 latency (ms)}, legend columns=3, legend style={at={(1.1,1.22)}, anchor=south}]
\xocurve{xoA}{*}{generated/eks-hpa/hpa-browse.csv}{A}\addlegendentry{A (monolith)}
\xocurve{xoB}{square*}{generated/eks-hpa/hpa-browse.csv}{B}\addlegendentry{B (4 services)}
\xocurve{xoC}{triangle*}{generated/eks-hpa/hpa-browse.csv}{C}\addlegendentry{C (8 services)}
\nextgroupplot[title={Checkout}, ymode=log, xmin=150, xmax=3200, xtick={200,400,800,1600}, xticklabels={200,400,800,1.6k}]
\xocurve{xoA}{*}{generated/eks-hpa/hpa-checkout.csv}{A}
\xocurve{xoB}{square*}{generated/eks-hpa/hpa-checkout.csv}{B}
\xocurve{xoC}{triangle*}{generated/eks-hpa/hpa-checkout.csv}{C}
\nextgroupplot[ymin=0, ymax=16, ylabel={CPU requested (cores)}, xmin=350, xmax=11000, xtick={500,1000,2000,4000,8000}, xticklabels={500,1k,2k,4k,8k}]
\addplot[gray, dashed, forget plot] coordinates {(100,\xo{hpa/cap}) (20000,\xo{hpa/cap})};
\foreach \c/\m in {A/*, B/square*, C/triangle*} {
  \edef\tmp{\noexpand\addplot[xo\c, mark=\m, mark size=1.4pt] table[x=rate, y=\c cap] {generated/eks-hpa/hpa-browse.csv};}\tmp
}
\nextgroupplot[ymin=0, ymax=16, xmin=150, xmax=3200, xtick={200,400,800,1600}, xticklabels={200,400,800,1.6k}]
\addplot[gray, dashed, forget plot] coordinates {(100,\xo{hpa/cap}) (20000,\xo{hpa/cap})};
\foreach \c/\m in {A/*, B/square*, C/triangle*} {
  \edef\tmp{\noexpand\addplot[xo\c, mark=\m, mark size=1.4pt] table[x=rate, y=\c cap] {generated/eks-hpa/hpa-checkout.csv};}\tmp
}
\end{groupplot}
\end{tikzpicture}
\caption{Autoscaling treatment in the cloud environment. Top: p99 latency in the recorded last \xo{hpa/window}\,s
of each \xo{hpa/step}\,s level (median across \xo{hpa/reps} repetitions, bars spanning the repetitions); a line
ends where a level was no longer sustained (more than 5\% errors, throughput below 90\% of the rate, or p99 above
\xo{hpa/slo}\,ms). Bottom: CPU requested by all pods of the configuration at the end of each level; the
dashed line is the shared quota of \xo{hpa/cap} cores.}
\label{fig:hpa}
\end{figure}

% TODO(results): text; sustained rates \xo{hpa/sustained/<wl>/<cfg>}, verdict \xo{hpa/verdict}.


\subsection{Autoscaling under a shared ceiling (RQ2, H2)}
\label{sec:res-hpa}

\begin{figure}[t]
\centering
\begin{tikzpicture}
\begin{groupplot}[group style={group size=2 by 2, horizontal sep=1.5cm, vertical sep=1.6cm}, xo, width=0.5\textwidth,
    height=4.6cm, xmode=log, xlabel={Offered load (requests/s)}]
\nextgroupplot[title={Browse}, ymode=log, xmin=350, xmax=11000, xtick={500,1000,2000,4000,8000}, xticklabels={500,1k,2k,4k,8k}, ylabel={p99 latency (ms)}, legend columns=3, legend style={at={(1.1,1.22)}, anchor=south}]
\xocurve{xoA}{*}{generated/eks-hpa/hpa-browse.csv}{A}\addlegendentry{A (monolith)}
\xocurve{xoB}{square*}{generated/eks-hpa/hpa-browse.csv}{B}\addlegendentry{B (4 services)}
\xocurve{xoC}{triangle*}{generated/eks-hpa/hpa-browse.csv}{C}\addlegendentry{C (8 services)}
\nextgroupplot[title={Checkout}, ymode=log, xmin=150, xmax=3200, xtick={200,400,800,1600}, xticklabels={200,400,800,1.6k}]
\xocurve{xoA}{*}{generated/eks-hpa/hpa-checkout.csv}{A}
\xocurve{xoB}{square*}{generated/eks-hpa/hpa-checkout.csv}{B}
\xocurve{xoC}{triangle*}{generated/eks-hpa/hpa-checkout.csv}{C}
\nextgroupplot[ymin=0, ymax=16, ylabel={CPU requested (cores)}, xmin=350, xmax=11000, xtick={500,1000,2000,4000,8000}, xticklabels={500,1k,2k,4k,8k}]
\addplot[gray, dashed, forget plot] coordinates {(100,\xo{hpa/cap}) (20000,\xo{hpa/cap})};
\foreach \c/\m in {A/*, B/square*, C/triangle*} {
  \edef\tmp{\noexpand\addplot[xo\c, mark=\m, mark size=1.4pt] table[x=rate, y=\c cap] {generated/eks-hpa/hpa-browse.csv};}\tmp
}
\nextgroupplot[ymin=0, ymax=16, xmin=150, xmax=3200, xtick={200,400,800,1600}, xticklabels={200,400,800,1.6k}]
\addplot[gray, dashed, forget plot] coordinates {(100,\xo{hpa/cap}) (20000,\xo{hpa/cap})};
\foreach \c/\m in {A/*, B/square*, C/triangle*} {
  \edef\tmp{\noexpand\addplot[xo\c, mark=\m, mark size=1.4pt] table[x=rate, y=\c cap] {generated/eks-hpa/hpa-checkout.csv};}\tmp
}
\end{groupplot}
\end{tikzpicture}
\caption{Autoscaling treatment in the cloud environment. Top: p99 latency in the recorded last \xo{hpa/window}\,s
of each \xo{hpa/step}\,s level (median across \xo{hpa/reps} repetitions, bars spanning the repetitions); a line
ends where a level was no longer sustained (more than 5\% errors, throughput below 90\% of the rate, or p99 above
\xo{hpa/slo}\,ms). Bottom: CPU requested by all pods of the configuration at the end of each level; the
dashed line is the shared quota of \xo{hpa/cap} cores.}
\label{fig:hpa}
\end{figure}

Figure~\ref{fig:hpa} shows the autoscaling treatment, \xo{hpa/cells} fresh deployments of A, B and~C. The monolith
sustained every level up to \xo{hpa/sustained/browse/A} browse and \xo{hpa/sustained/checkout/A} checkout requests
per second, and it did so with a single replica. At those rates its application container stayed below the
autoscaler's 70\% target, so it requested only its starting \xo{hpa/topcpu/browse/A} cores, 3.5 for the application
and 4 for the database. The decomposed configurations sustained every level only up to
\xo{hpa/sustained/browse/B} (B) and \xo{hpa/sustained/browse/C} (C) browse requests per second and
\xo{hpa/sustained/checkout/B} and \xo{hpa/sustained/checkout/C} checkout requests per second. Their autoscalers did
act: beyond these rates they added replicas until the configurations requested the whole ceiling of
\xo{hpa/cap} cores. The new replicas joined cold while the load kept rising, however, and p99 broke the
\xo{hpa/slo}\,ms target in at least one repetition. B reached the monolith's highest sustained browse level,
\xo{hpa/top/browse/B} requests per second, in every repetition only at the ceiling, with \xo{hpa/topcpu/browse/B}
cores requested against the monolith's \xo{hpa/topcpu/browse/A}. At the first level every configuration scaled up
briefly, because the metrics pipeline still reported the CPU of the warm-up that had just ended; the next level
undid it, and the 30\,s idle period before the ramp was evidently too short to prevent it.

\textbf{H2 is therefore falsified under autoscaling as well.} One detail of the rule in
Section~\ref{sec:hypotheses} needed a margin. At low rates, where no autoscaler had acted, B and~C requested a few
millicores less than the monolith because their static shares are rounded down. Read literally, the rule would
count that as sustaining a rate with less CPU. We therefore required the difference to exceed 1\% of the ceiling.
This margin was fixed after we saw the case, and with it no decomposed configuration sustained a rate the monolith
sustained while requesting less CPU.


\subsection{Summary of the hypothesis tests}

\begin{table}[t]
\caption{Hypotheses, falsification conditions and outcomes in the local environment.}
\label{tab:verdicts}
\centering\small
\begin{tabular}{@{}lp{0.52\linewidth}l@{}}
\toprule
& Falsified if & Outcome \\
\midrule
H1 & any microservice configuration has a significantly lower p99 than A at the lowest rate & \xo{verdict/H1} \\
H2 & A matches or exceeds every microservice configuration at every concurrency level, or no $\beta < \beta_A$ & \xo{verdict/H2} \\
H3 & mesh $<25\%$ of end-to-end latency at six hops, or per-hop cost shrinking with depth & \xo{verdict/H3} \\
\bottomrule
\end{tabular}
\end{table}

Table~\ref{tab:verdicts} summarises the three tests. Decomposition costs latency at low load (H1) and the mesh's
cost compounds with depth (H3), as predicted. The central hypothesis, that horizontal decomposition eventually
overtakes the monolith (H2), is falsified for an equal CPU budget: the monolith's capacity was four to sixteen times
that of the decomposed configurations, and no crossover appeared at any load.

